% \newpage
% \begin{center}
%   \textbf{\large АННОТАЦИЯ}
% \end{center}


% Объектом исследования данной работы стали системы частиц, взаимодействующих посредством обобщенного потенциала Леннарда-Джонса с переменной степенью притяжения.
% С помощью данных систем выявлялось влияние дальнодействия притяжения на фазовые диаграммы, роль дальнодействия в транспортных свойствах, а также влияние на спектры возбуждений.
% Для решения поставленной задачи методом молекулярной динамики были смоделированы системы.
% Пост обработка проводилась с использованием MATLAB и python.

% В ходе работы были проведены моделирования систем с различным дальнодействием притягивающей ветви потенциала взаимодействия.
% Рассчитаны фазовые диаграммы систем, их транспортные свойства и спектры возбуждений на бинодали жидкость-газ.
% Предложен новый метод классификации частиц на конденсат, газ и поверхность в системах с фазовым расслоением.
% На основании полученных результатов, было впервые выявлено влияние дальнодействия притяжения в молекулярных системах на фазовые диаграммы, положения тройных и критических точек и корреляции транспортных свойств со спектрами возбуждений на бинодали жидкость-газ.
% С помощью нового метода классификации частиц, основанного на методе кластеризации DBSCAN (Density-Based Spatial Clustering of Applications with Noise), был разработан метод построения фазовых диаграмм и анализа нуклеации.

% \onehalfspacing
% \setcounter{page}{2}

\newpage
\onehalfspacing
\setcounter{page}{2}
\renewcommand{\contentsname}{\centerline{\large СОДЕРЖАНИЕ}}
\tableofcontents

\newpage
\begin{center}
  \textbf{\large ВВЕДЕНИЕ}
\end{center}
\addcontentsline{toc}{chapter}{ВВЕДЕНИЕ}

В современном мире практически не возможно обойтись без информационных технологий. Цифровизации подвергается каждая сфера жизни человека: социальные сети, игры, маркет-плейсы и нейросети: без этого уже сложно представить нашу жизнь. С ростом цифрового следа пользователя, ростёт и количество ценных информационных активов, которые он хранит в сети «Интернет»: цифровые финансовые активы, криптовалюты, ценные игровые артифакты и многое другое. Если процесс передачи материальных активов по наследству уже давно отлажен, то с цифровыми активами мы сталкнулись недавно и ещё имеем недостаточно опыта для работы с ними.
\textbf{Актуальность} работы обусловлена потребностью системно управлять цифровыми активами: Хранить, систематизировать и передавать их по наследству или в пользу государства.
\textbf{Целью выпускной квалификационной работы} является разработка \\ требований для государственной информационной системы для управления, каталогизации, систематизации цифровых активов граждан Российской Федерации.

Для достижения поставленной цели необходимо решить следующие \textbf{задачи}:
\begin{enumerate}
\item Анализ существующих подходов и нормативно-правовой базы в области менеджмента цифровых активов, как на государственном уровне, так и коммерческих решений. 
\item Исследовать классификацию и жизненный цикл цифровых активов, требования к их хранению и обработке.
\item Составить функциональные и нефункциональные требования к государственной информационной системе управления цифровыми активами и разработать концептуальную архитектуру системы, а также предложить модель интеграции с в существующуюю правовую систему и другие институты, такие как нотариат. 
\end{enumerate}

\textbf{Объект исследования:} цифровые активы, как часть мировой экономической системы.

\textbf{Предмет исследования:} совокупность требований к системе управления цифровыми активами, учитывающая текущую нормативно-правовую базу, уровень развития технологий  общества.

Данная выпускная квалификационная работа состоит из введения, трех глав, заключения и списка использованной литературы. В первой главе рассматривается a, Во второй главе b, В третьей главе c. В данной работе представлены: w таблиц, x рисунков и y формул, а также z приложений.